% ===== Envoi: how to think of that rose =====
\section*{Envoi: How to Think of That Rose}
\addcontentsline{toc}{section}{Envoi: How to Think of That Rose}
\label{sec:envoi}

\noindent
This is not a conclusion---the paper has said why it can have none---but a returning, of the kind the paper has been about: a coming back to where it began, the rose in the garden, carrying home what the going has gathered. So let the last word not be an argument but a looking-again at that one flower, and at how we might think of it.

The remark that began this paper grieved that the rose, however hard it labours, blooms but a week. And it is true, and it cannot be otherwise: that the flower will fall is the natural turning of the cycle, and no tending we could devise will change it. To fix our care upon the falling---to set ourselves against it, to water and fuss and clutch in the hope of holding the bloom past its hour---is the surest way to rot the root; it is to grasp, and the one who grasps loses. So let our attention fall elsewhere than on the falling. \emb{Rather than attend to the rose's withering, which we cannot change, let us attend to how the rose was generated---for generativity, unlike the bloom, can perpetuate itself.}

If it is a rose in the garden, then: tend its root. The flower is the part that is seen and the part that falls; the root is the part that is unseen and the part that goes on. To care for the root is not to neglect the flower but to love it rightly---to love it in the only way that brings, after this bloom, another. This is what it has meant, all through these pages, not to forget the origin: to keep faith with the root from which the visible flowering came, and to which, when it falls, the flower returns.

And if it is a rose given to one you love, then do not forget how this rose came to be in your hand, nor why it is here. It is not, first of all, a thing; it is the showing of a love, and the love is itself a generativity---a generativity that, kept and not clutched, returned to and not grasped at, will bring you, again and again, many roses more. The single flower in the hand will fall, as the one in the garden falls; but the love it shows, if it is let to be the spiral and not made into the held and closing circle, does not fall, and from it the roses come and come. To receive the rose rightly is to receive, through it, the generativity it bears witness to---and to know that what one holds is not the flower, which is mortal, but the root of a love, which goes on giving flowers.

So we need not save the rose. We could not, and the trying would be the rotting. We need only keep faith with the generativity that made it---tend the root, remember the love---and the roses will keep coming, each brief, each whole, each falling in its turn to feed the next. That is the sustainability this paper has been about, said now in the plainest way it knows: not a flower that does not fall, but a love that does not stop generating flowers. Let the bloom be brief. Tend the root. The roses will come.
